\documentclass[11pt]{article}
\usepackage[final]{acl}
\usepackage{times}
\usepackage{latexsym}
\usepackage[T1]{fontenc}
\usepackage[utf8]{inputenc}
\usepackage{microtype}
\usepackage{inconsolata}
\usepackage{graphicx} \title{Critical Review of ``Deep Learning on Object-Centric 3D Neural Fields'' (T-PAMI 2024)}
\author{First Last \
Department of Computer Science, University Name \
\texttt{email@domain.edu}}
\date{} \begin{document}
\maketitle \begin{abstract}
This review examines the paper “Deep Learning on Object-Centric 3D Neural Fields” (Ramirez et al., IEEE T-PAMI 2024) from the perspective of a computer vision graduate student. We summarize the proposed nf2vec framework, which aims to embed implicit 3D neural field representations into compact vectors for downstream deep learning tasks. The promise of unified neural representations for 3D data is highlighted, along with the initially exciting prospects of nf2vec enabling a single pipeline for diverse 3D tasks (classification, generation, interpolation, etc.). We then provide a critical evaluation of the approach, discussing its strengths (scalability, elegance, modality-generalization) and weaknesses (reliance on shared network initialization, entanglement of geometry/appearance, limited data augmentation). A brief outline of a seminar presentation on the paper is included, and a placeholder notes the omission of additional experimental attachments due to time constraints. The tone is constructive yet skeptical, reflecting on how initial enthusiasm was tempered by deeper understanding of nf2vec’s practical challenges.
\end{abstract} \section{Introduction}
Understanding the 3D world has long been a core goal of computer vision, but it has been impeded by fragmentation in how 3D data is represented. Traditional representations—voxel grids, point clouds, polygon meshes—each require specialized algorithms and neural network architectures, with no single standard for storing and processing 3D information. This lack of a unified representation has led to isolated pipelines for each modality and difficulties in transferring methods across them. Discrete 3D representations also suffer from trade-offs like fixed resolution and high memory usage. In recent years, Neural Fields (NFs) have emerged as an appealing alternative to discrete geometry. An NF represents a signal (such as the shape or appearance of a 3D object) as a continuous function defined over $\mathbb{R}^3$, typically parameterized by a coordinate-based neural network. For 3D objects, common instantiations include implicit surface fields like signed/unsigned distance functions (SDF/UDF) and occupancy fields (OF) which define geometry, as well as radiance fields (RF) which encode both geometry and appearance (as in Neural Radiance Fields, NeRF). Neural fields decouple memory from resolution: a fixed-size MLP can represent details at arbitrary precision, unlike voxels or images. Moreover, the same neural network architecture can learn different field types, hinting at a unified framework for representing 3D objects. This potential to standardize 3D representation has sparked great interest. Researchers speculate that NFs could become a standard format for 3D assets in the near future, with databases of digital objects stored as networks instead of meshes. Such a paradigm shift would demand new ways to process NFs directly in learning pipelines. A key research question thus arises: Can we feed 3D neural fields directly into deep learning models for downstream tasks, instead of converting them to point clouds, voxel grids, or images? For example, could one classify an object by operating on its NeRF representation without ever rendering an image? Initially, this prospect was very exciting — it promised to eliminate conversion overhead and unify 3D understanding under one paradigm. However, NFs are essentially neural networks (with their information encoded in weights), so traditional deep learning methods cannot directly treat them as input in the same way as tensors of pixels or points. Early attempts to process NFs in deep learning pipelines revealed challenges. Functa (Dupont et al., 2022) explored encoding an entire dataset of NFs into a single shared neural network conditioned on per-sample codes. In theory this provides latent descriptors (the conditioning vectors) that could be used for tasks, but in practice the single network struggled to fit all samples faithfully. Representing many objects with one set of weights proved too difficult, yielding suboptimal representations. By contrast, other works showed that fitting one neural network per sample (each object has its own NF) achieved much higher fidelity. This individual-network approach aligns better with real-world use (each 3D object can be independently parameterized) and avoids forcing one model to be “one size fits all.” The question remained how to leverage those individually trained NFs for tasks like classification or generation in a deep learning framework. The paper “Deep Learning on Object-Centric 3D Neural Fields” addresses this question with a solution called nf2vec. The idea of nf2vec — encoding each neural field into a compact vector in a single forward pass — was initially very compelling. It promised a unified representation for any 3D object stored as an NF, enabling us to plug these embeddings into standard machine learning models. If successful, this would allow 3D tasks to be approached with the same ease as image or text tasks, using a common feature space regardless of whether the input was an SDF, an OF, or even a NeRF. The prospect of training one encoder and then seamlessly applying classical pipelines (classifiers, generative models, etc.) on the resulting latent codes generated considerable enthusiasm. In theory, nf2vec could bridge the gap between the continuous implicit world of NFs and the discrete tensor world of deep learning algorithms. This review will first summarize the nf2vec framework and its contributions. We then critically evaluate how well it fulfills its promise, discussing both the attractive aspects (e.g. elegance and generality of the approach, experimental successes) and the drawbacks or technical caveats (e.g. conditions required for it to work, limitations in performance and flexibility). In hindsight, the initial excitement around nf2vec’s unification of 3D representations is tempered by a nuanced understanding of its assumptions and limitations. We articulate these insights in the context of scalability, representational power, and practicality. Additionally, we include an outline of a seminar presentation delivered on this paper, which mirrors how the material can be conveyed in an academic setting. Finally, we acknowledge omitted experimental details and attachments, noting that additional exploration was beyond the scope of this review. \section{Summary of the Paper}
\subsection{The nf2vec Framework}
At a high level, nf2vec provides a way to convert a learned neural field into a fixed-length vector embedding in one inference step. This embedding is intended to capture the essence of the 3D object represented by the NF (its geometry, and optionally appearance) in a task-agnostic manner. Once an NF is distilled into this latent code, it can be fed into typical deep learning machinery for various downstream tasks just as one would use image feature vectors or word embeddings. Crucially, the pipeline avoids explicitly sampling the NF on a grid or point set; instead, it “looks” at the NF in its native form (as a set of network weights) to produce the encoding. The nf2vec framework consists of two primary components: an encoder network that produces the latent vector from the NF, and an implicit decoder network that can use the latent vector to reconstruct the original field’s output. Figure 1 of the paper provides an overview: the left side illustrates that neural fields (NFs) offer a unified representation of 3D data, the center shows nf2vec compressing an NF into a code vector, and the right side depicts how these codes can interface with common deep learning tasks (classification, generation, retrieval, etc.). In essence, nf2vec is a form of autoencoder operating on neural field parameters rather than on images or point clouds. Encoder Architecture: A naive way to encode an NF would be to feed its entire weight vector into another very large network, but this would be infeasible for NFs with hundreds of thousands of parameters. The authors design a more efficient encoder that exploits the layered structure of MLP-based NFs. They note that a typical NF (for example, an MLP with $L$ hidden layers of width $H$) has weight matrices for each layer; if you flattened all weights, an enormous encoder (on the order of the number of NF parameters times the embedding size) would be required. Instead, nf2vec stacks the NF’s weight matrices and bias vectors into a big matrix $P$ of shape $S \times H$ (where $S$ is the total number of weight “rows” and $H$ is the common hidden layer size). Each row of $P$ corresponds to one neuron’s weights (and bias) from one layer, after zero-padding where necessary to align dimensions. The encoder then applies a series of linear layers with batch normalization and ReLU to each row of this matrix independently, transforming each row into a new row-feature of the same length. Finally, a max-pooling aggregation is performed across the row dimension, collapsing the matrix into a single 1×$H$ embedding vector. This design (illustrated in Fig. 2 of the paper) ensures the encoder’s complexity grows roughly linearly with the number of weights, instead of quadratically. By processing weight rows independently with shared transform layers, the encoder can scale to large NFs without a blow-up in its own parameter count. Table I in the paper quantitatively confirms that their encoder has orders of magnitude fewer parameters than a hypothetical dense encoder for typical NF sizes. One important detail is that nf2vec assumes all input NFs share a common architecture and even a common weight ordering. The authors enforce that each NF is trained from a shared random initialization so that corresponding weight rows play similar roles across different NFs. Without this, two networks representing the same shape could have completely different weight patterns (due to symmetric permutations of neurons or different random seeds), making it impossible for the encoder to interpret rows consistently. Adopting a fixed initialization aligns the weight space across samples and was found to be crucial for nf2vec to converge and produce meaningful embeddings. (This constraint is an important point we will revisit in the critical evaluation.) Decoder and Training: The decoder in nf2vec is another neural network that takes as input the learned embedding (the output of the encoder) and a 3D coordinate $p=(x,y,z)$, and outputs the field value $\hat{f}(p)$ at that point. In other words, the decoder is an implicit function network (similar in form to the original NF) conditioned on the code. The decoder architecture largely mirrors a standard MLP used for the original NF, augmented with the ability to accept the latent code. In practice, the embedding is concatenated with the coordinate (and injected at multiple layers via skip connections in the implementation) so that the decoder can modulate the output based on the particular object’s code. The goal of training nf2vec is for the decoder to reproduce the behavior of the original NF \emph{without directly copying its weights} – i.e., to reconstruct the shape or radiance field from the code alone. The encoder and decoder are trained together in an unsupervised fashion using a large set of known neural fields (for example, many NFs representing shapes from ShapeNet or similar datasets). For each training NF, the encoder produces an embedding, and then the decoder is asked to predict the NF’s output for various sample coordinates. The training loss penalizes the difference between the decoder’s prediction $\hat{f}(p)$ and the original NF’s true value $f(p)$ across a set of sample points $p$ drawn from the object’s domain. Essentially, this encourages the embedding to be a sufficient summary of the NF such that the decoder can imitate the NF’s mapping. Notably, the decoder is not trying to reconstruct the NF’s weights, but rather to match its input-output function. This focus on function reconstruction means the latent space is tuned to preserve the object’s geometric (and visual) features, not the exact parameter values of the original NF. After training, the encoder is “frozen” and can be used to embed new, unseen NFs in a single forward pass. If needed, the trained decoder can also be used to convert an embedding back into a discrete representation (e.g., by sampling many points to produce a point cloud or by volume rendering to produce images), although for many tasks one might not require this step. The authors demonstrate nf2vec on several types of neural fields:
For pure geometry fields like SDF, UDF, and OF, the decoder is a scalar function network that outputs a single value (distance or occupancy) for a 3D coordinate. Training these is straightforward since ground-truth values (e.g., signed distances) can be obtained from the original shape data to supervise the NF and thus to supervise nf2vec’s decoder. In their experiments, NFs for shapes were often obtained by fitting SDFs to mesh data or UDFs to point clouds; nf2vec was then trained on those.
For Neural Radiance Fields (NeRFs), which encode both shape and appearance, the decoder must output both color and density (often represented as an RGB+$\sigma$ output). The paper extends nf2vec to handle NeRFs by using a decoder with a 4-channel output (3 color channels and 1 density channel) and training it with a differentiable volumetric rendering process. Specifically, they render 2D images from the decoder’s output and compare them to ground-truth views of the original NeRF, using an $L_1$ loss on pixel colors. They give higher weight to foreground (object) pixels and lower weight to background in the loss to encourage the latent to capture object-specific content. This way, the embedding learns to encode both geometry and appearance in a single vector, and the decoder learns to reproduce a radiance field whose renders match the input NeRF’s renders. Despite the complexity of jointly representing shape and color, nf2vec is shown to handle NeRF inputs as well, using the same encoder architecture (the difference is only in the decoder’s output head and the training loss).
\subsection{Experiments and Results}
After training nf2vec on a given domain of NFs, one can convert all those neural fields into their latent vector form. The paper evaluates these embeddings on a variety of downstream tasks to demonstrate their usefulness. Importantly, all tasks are performed without ever explicitly converting the NFs to meshes, point clouds, or images; the pipeline operates entirely in the space of learned embeddings, showcasing a new way to do “deep learning on 3D neural fields” directly. Shape Classification: The authors tackle 3D object classification by feeding nf2vec embeddings into a simple feed-forward classifier (a 3-layer perceptron). They consider multiple benchmark datasets originally in different 3D formats (e.g., ModelNet40 for CAD models, ShapeNet, etc.), convert all shapes to NFs (SDFs, UDFs, or OFs depending on the case), embed them with nf2vec, and then train/test a classifier on the resulting codes. Remarkably, this approach achieves classification accuracies that are competitive with specialized 3D neural networks designed for each input modality. For example, on ModelNet40 (where shapes were represented as UDFs of point clouds), nf2vec reached around 87% accuracy, comparable to state-of-the-art point cloud classifiers. On a mesh-based dataset (Manifold40), it obtained about 86%, and on a ShapeNet10 subset (voxelized objects), about 93%. These numbers are only marginally below dedicated architectures like PointNet++ or voxel CNNs, yet nf2vec used one unified pipeline for all, with the same classifier architecture applied to all embedding types. This underscores one major advantage of nf2vec: modality-independence. Competing methods had to be tuned to point clouds or volumes specifically, whereas nf2vec’s embeddings serve as a common denominator. Additionally, inference time is improved — classifying an object via its nf2vec code is faster than converting an NF to, say, a $2048$-point cloud and then running a point network. The paper reports that classifying a shape from its NF took only about 1ms with their approach, whereas even a single image-based classification (rendering the NeRF to an image and using ResNet) could take nearly 100ms. Thus, nf2vec offers efficiency gains by operating directly on the implicit representation. They also experimented with part segmentation (predicting labels for each part of a 3D object) using nf2vec. In that scenario, a global embedding alone cannot produce per-point labels, so they combined the nf2vec code with point-level features in a segmentation network. The results showed that even without explicitly encoding spatial details, the nf2vec embedding provides useful object-level context for segmentation, achieving respectable part segmentation performance. This was an interesting demonstration that the learned latent captures not just class identity but some information about the object’s finer structures. Reconstruction and Retrieval: Because nf2vec is essentially an autoencoder, one measure of representation quality is whether the decoder can reconstruct the original shape from the embedding. Qualitatively, the paper’s Fig. 4 and others show that shapes reconstructed via the decoder from the compressed code look very similar to those obtained by directly rendering the original NF. Even though the embedding might be only 1024 numbers instead of the hundreds of thousands of weights in the NF, the important geometric features are preserved. They also visualize the latent space structure using t-SNE: embeddings naturally cluster according to object category (e.g., chairs together, tables together) despite the model not being explicitly trained for classification. This suggests the nf2vec latent space is semantically meaningful. In a shape retrieval task (finding nearest neighbors in code space), the model indeed retrieves objects of the same class, indicating that similarity in embedding space correlates with similarity in shape. Interpolation: One hallmark of a well-structured latent space is that interpolating between embeddings yields plausible intermediate representations. The authors provide results for interpolating between two shape embeddings and between two NeRF embeddings. Linearly blending the latent vectors produces a smooth morphing of geometry; for instance, an airplane gradually turns into a rifle in the decoded shape visualization (Fig. 6 of the paper). In the case of NeRF (which includes appearance), interpolation causes both shape and color to transition smoothly. The intermediate latent produces a 3D object that looks like a meaningful hybrid of the two endpoints, and when rendered from different viewpoints it remains a coherent object (demonstrating 3D consistency of the interpolation). This is a notable result: averaging two sets of network weights directly would typically yield something meaningless (e.g., blobby or broken shapes), but averaging nf2vec codes yields a sensible object. The authors point out that direct averaging of two NeRF weight sets produced blurred outputs lacking structure, whereas averaging in latent space gave a sharper result. This illustrates that nf2vec learned a latent representation where linear combinations correspond to logical combinations of object properties, a desirable property for generative modeling. Generation and Cross-Modal Mapping: Going further, the paper explores generative modeling in the nf2vec latent space. They train a simple GAN (Generative Adversarial Network) on the embeddings of shapes, which allows sampling new random codes and decoding them to novel 3D shapes. The shapes generated by this latent-GAN exhibit realism and diversity, and because the decoder is an implicit field network, these shapes can be rendered at arbitrary resolution or converted into any representation as needed. In addition, nf2vec enables some unique cross-representation tasks. For example, they train a mapping network that takes as input an embedding from a geometry-only field (like a UDF of an object) and predicts the embedding of that same object’s NeRF. This effectively adds appearance to a pure geometry representation by going through the latent domain. In their results (Fig. 21 of the paper), a point cloud-derived UDF embedding is converted into a NeRF embedding; decoding that yields a colored 3D model, showing the method can transfer an object from one NF representation to another via learned latent mappings. They also demonstrate shape completion: mapping an embedding of an incomplete object (partial shape) to the embedding of a complete shape, thus using latent space to perform completion tasks. These examples underscore nf2vec’s generality: once objects are in the vector space, standard machine learning models (GANs, MLP mappers) can operate on them to achieve tasks like generation, transformation, and translation between different NF modalities. Finally, the paper provides a comparison with other contemporary approaches for learning on neural fields. They compare nf2vec to methods like DWSNet, NFN, and NFT (which incorporate weight-space symmetries into their architectures), as well as to the shared-network approach (Functa, DeepSDF). The results show that nf2vec achieves the best balance of reconstruction fidelity and task accuracy. In particular, methods relying on a single shared network for all data (Functa, DeepSDF) tend to sacrifice reconstruction quality and can struggle with diversity. Per-field methods like nf2vec and its peers maintain high fidelity; among these, nf2vec stands out for its strong classification performance (outperforming or matching others on NF classification benchmarks). The authors highlight that they established the first benchmark for neural field classification, and nf2vec achieves state-of-the-art results on it (about on par with Functa in accuracy, but without Functa’s limitations, and with far better shape fidelity). These comparisons reinforce that nf2vec is a leading method for directly processing implicit 3D representations, pushing the field toward treating NFs as first-class input to neural networks. \section{Critical Evaluation}
The nf2vec approach marks an important step toward unifying 3D representation learning. It elegantly demonstrates that implicit neural fields can be integrated into deep learning workflows without explicit meshing or rasterization. However, a closer examination reveals several caveats and challenges. In this section, we discuss the pros and cons of nf2vec in detail, reflecting on how our initial excitement has matured into a more critical understanding. \subsection*{Pros}
Unified Representation & Modality-Generalization: One of nf2vec’s greatest strengths is that it provides a common embedding space for heterogeneous 3D representations. It successfully embeds signed distance fields, unsigned distance fields, occupancy grids, and radiance fields under the same framework. Consequently, the same downstream network (e.g., a classifier or regressor) can be applied across all these input types, which was previously impossible with specialized models. This generality is elegant – nf2vec effectively acts as a “universal adapter” for 3D data, enabling a unified 3D deep learning pipeline for shapes that were originally in incompatible forms.
Task-Agnostic Embedding (Versatility): The latent code produced by nf2vec is not tailored to any single task; yet it proves effective for a wide range of applications including classification, retrieval, part segmentation, shape interpolation, and even generative modeling. This is a positive indication that the embedding captures semantically meaningful and comprehensive information about the object. It is impressive that a single 1024-dimensional vector (for instance) can encapsulate enough detail to allow both fine-grained discrimination (distinguishing object categories, aiding segmentation) and faithful reconstruction of the object’s shape/appearance. In other words, nf2vec finds a sweet spot in representation learning – compressing the NF while preserving essential content. The fact that interpolating between embeddings yields smooth transitions and that embeddings cluster by object category without supervision further attests to the meaningfulness of the learned space.
Scalability and Efficiency: The encoder design is cleverly optimized for scale. By operating on each weight-row independently and using pooling, the encoder’s parameter count grows much more slowly than the size of the input NF. This means nf2vec can in principle handle very large NFs (with more layers or neurons) without becoming impractical to train or run. Additionally, inference with nf2vec is efficient: embedding an NF is a single forward pass, and once embedded, tasks like classification are extremely fast (on the order of milliseconds). The authors demonstrated that direct NF classification via nf2vec was two orders of magnitude faster than pipelines that require rendering images or sampling point clouds. Real-time or on-device 3D understanding could benefit from such speed-ups. The approach also supports batch processing — multiple NFs can be encoded in parallel — which is beneficial for throughput in practical applications.
Competitive Performance: Considering its generality, nf2vec achieves admirable performance on specialized benchmarks. Its shape classification accuracies on ModelNet40, ShapeNet10, and others are only slightly below state-of-the-art methods that are specifically tuned for those data formats. This is a notable accomplishment: nf2vec did not have the luxury of exploiting modality-specific cues (e.g., learning directly on point sets or meshes), yet it nearly matches those methods while using a uniform approach. Moreover, nf2vec outperforms prior neural-field-specific methods on the new NF classification benchmark the authors propose. It also yields high-fidelity reconstructions — the decoder’s outputs are almost as accurate as the original NFs, indicating minimal information loss in the embedding. In qualitative terms, one can hardly distinguish a surface or image produced from the compressed code versus the ground truth NF, which speaks to the efficacy of the learning scheme.
Smooth Latent Space & Generative Potential: The learned latent space appears well-behaved and smooth. We observe that linearly interpolating between two embeddings produces intermediate shapes that are plausible and maintain 3D coherence. This property is non-trivial and suggests nf2vec did not overfit to individual objects but rather organized the space in a meaningful way. The ability to perform arithmetic in latent space (e.g., for shape morphing or attribute transfer) opens up new possibilities for creativity and data augmentation. Furthermore, since standard generative models (GANs, diffusion models) can be trained on these embeddings, nf2vec effectively enables generative modeling of 3D objects in the neural field domain. This is exciting because it leverages the benefits of implicit fields (high-quality continuous representations) while operating in a compact vector space. The success of nf2vec in cross-modal mappings (e.g., geometry-only to NeRF embedding) also hints that the embedding space may serve as a lingua franca for different types of 3D data, which is a conceptually beautiful result.
Technical Elegance: From an architecture perspective, nf2vec’s design is elegant and simple. It does not rely on complex architectural tricks beyond the row-wise encoding and pooling. There is a certain simplicity in treating the NF weights almost like a “set” or unordered list of features and using a symmetric function (max-pool) to gather them. This is reminiscent of deep set architectures (like PointNet for point clouds), drawing a parallel between unordered point sets in geometry and unordered weight sets in an MLP. The approach sidesteps the need for heavy 3D convolutions or sampling procedures, which are often bottlenecks in 3D deep learning. Instead, it directly consumes the MLP parameters, which is a novel viewpoint. The method’s ability to leverage existing deep learning techniques (classifiers, etc.) on top of these embeddings is an attractive quality for integration into broader systems.
\subsection*{Cons and Limitations}
Despite its merits, nf2vec has several limitations and assumptions that became apparent upon closer analysis:
Reliance on Shared Initialization: A notable constraint is that all NFs must be trained with an identical random initialization (and architecture) in order for the encoder to function properly. The encoder essentially assumes that the $i$-th row in matrix $P$ for one NF is “comparable” to the $i$-th row in another NF. Without a consistent initialization, this assumption breaks due to weight space symmetries (e.g., permuting hidden neurons can drastically change the weight matrix while representing the same function). The authors themselves acknowledge that MLP weights have symmetries (neurons can swap with no effect on output) and that nf2vec must contend with this issue. They adopt the strategy of identical initialization as a practical fix to align weights across objects. However, this requirement can be seen as a rigidity: in scenarios where NFs come from different sources or are trained independently with different inits or architectures, nf2vec cannot directly embed them into a common space without retraining or finetuning. It also ties nf2vec to the specific MLP architecture used during training – any deviation (say, an NF with a different number of layers or neurons) would violate the assumptions. In summary, nf2vec trades some flexibility for representational consistency. This might hinder its adoption if one cannot ensure all implicit models are trained with the same template.
No Straightforward Online Augmentation: Data augmentation is a cornerstone of robust deep learning models, especially in vision (random rotations, scaling, jitter, etc. for 3D objects). In nf2vec’s case, performing augmentations directly on neural fields or their weights is non-trivial. One cannot easily “rotate an NF” by just rotating its weight matrix; one would have to either rotate the input coordinates (which would require re-fitting the NF or applying a coordinate transform trick) or manipulate weights in a complex way. The authors mention that they resorted to augmenting the training dataset offline – for instance, by generating new NFs from transformed versions of shapes – rather than doing on-the-fly augmentation. This is a limitation because it reduces the diversity seen during training and complicates the pipeline (you need to prepare many NFs in different poses a priori). Some recent works (they cite E-StitchUp) try to augment at the embedding level, but generally this is an unsolved problem. The inability to perform simple augmentations like random perturbations or rotations directly on the NF weights means nf2vec might struggle with invariances or require more training data to compensate. This stands in contrast to, say, point cloud methods where a random rotation can be applied to the input with ease during training.
Entangled Geometry and Appearance: In nf2vec’s handling of NeRFs, a single embedding vector encodes both the object’s shape and its visual appearance. The paper explicitly notes that “appearance and geometry are intrinsically entangled in our framework” when processing radiance fields. This entanglement has a downside: it is impossible to adjust an object’s color/texture without affecting its shape via the embedding (or vice versa). For some applications, one might desire separate control (e.g., keep the geometry but change the material/texture). Because nf2vec treats the entire NeRF as one monolithic chunk of weights to embed, all aspects are fused. In contrast, some other approaches to generative NeRFs factor the latent space into shape and appearance codes explicitly. Thus, nf2vec sacrifices this flexibility. Moreover, the single-code approach could make the embedding space more complex to learn, since it must account for variations in both geometry and color simultaneously. In the current state, nf2vec is better suited to tasks where geometry and appearance are not required to be disentangled.
Only Object-Centric (Not Scene-Level): By design, nf2vec processes one object’s NF at a time. It assumes a clean object-centric NF (often with the object near the origin, isolated). Extending this to full scenes with multiple objects or large environments is not straightforward. A scene-level implicit representation (like a large NeRF of a room or a continuous SDF of an environment) could in principle be treated as an NF, but embedding an entire scene might require a much larger network or code, and it’s unclear if the current encoder would scale or if the latent space would meaningfully encode multi-object structure. The authors hint at this in their discussion of future work: handling scenes with multiple interacting NFs would be a next step. This is not so much a flaw of nf2vec as it is an scope limitation, but it means that the method is currently confined to single objects. Applications like relational scene understanding or multi-object classification would need additional developments (perhaps combining multiple embeddings or designing an encoder that can aggregate multiple NFs).
Performance Trade-offs: While nf2vec achieves high accuracy in many cases, it occasionally lags behind specialized methods. For instance, the very best point cloud classifiers slightly outperform nf2vec on shape benchmarks, and certain tailored models might do better on specific tasks (e.g., a network specifically designed for segmentation might outperform using nf2vec embeddings plus a generic segmentation head). This is expected because nf2vec’s strength is in its generality. Nonetheless, it’s a trade-off worth noting: users of nf2vec pay a small penalty in raw performance for the benefit of a unified approach. In critical applications where every percentage of accuracy matters, this could be a concern. Additionally, the embeddings are of fixed size (1024 in the paper) regardless of the complexity of the object. This fixed bottleneck might struggle to capture extremely intricate shapes or scenes with fine details, whereas an explicit representation could always add more points or voxels to cover detail. In their experiments, nf2vec reconstruction was very good but not perfect – some very fine details or high-frequency texture might be slightly smoothed by the compression. Thus, there is an inherent information bottleneck in any autoencoder, including nf2vec.
Implicit Assumptions and Training Complexity: Training nf2vec itself is not a trivial task. It requires a large dataset of neural fields (which in turn may require fitting many NFs to raw data). The process of generating those NFs (with consistent initialization) and then training a sizable encoder-decoder on them is quite involved. There are also subtle assumptions: for example, the input and output layers of SIREN-based NFs were discarded before encoding, as noted in the seminar slides (because sine-based input layers might not carry object-specific info). Such domain-specific adjustments need to be done carefully. The method also assumes that the NF training procedure yields networks that are similar enough for the encoder to handle. If one NF was poorly trained or overfit, its weights might be out-of-distribution for the encoder. In practice, the authors navigated these issues well, but they highlight that nf2vec is a complex pipeline with many stages (data preparation, NF fitting, representation learning) which could make it harder to use or reproduce than more direct end-to-end models on raw data.
In summary, nf2vec is an innovative framework that demonstrates the feasibility of deep learning directly on implicit representations. It brings clear advantages in unification and flexibility of output, but it introduces its own prerequisites (shared NF format and training regime) and does not fully escape the challenges inherent to weight-space learning (such as symmetry issues and interpretation of neural weights). The initial enthusiasm for a “one-size-fits-all” 3D representation is warranted, but tempered by recognition that nf2vec works under certain conditions that may limit its universality. As users, we must be mindful that the approach may need further refinement (e.g., methods to relax the shared initialization requirement or to disentangle latent factors) before it can become a drop-in solution for all 3D learning problems. The authors are optimistic and so are we, but our optimism is coupled with a realistic understanding of the current limitations. \section{Seminar Presentation Outline}
In a seminar presentation accompanying this paper, the major points and flow of “Deep Learning on Object-Centric 3D Neural Fields” were structured as follows (slide titles in bold, with key contents described):
Motivation – The 3D Representation Problem: Introduction to the challenges of 3D data representation. Emphasized how 3D data has been fragmented into voxels, point clouds, meshes, etc., each requiring specialized processing. Highlighted that no single standard or pipeline exists for all 3D data, and discussed limitations of discrete representations (e.g. fixed resolution, high memory).
Neural Fields as Continuous Representations: Introduced Neural Fields (NFs) as a new paradigm for representing 3D content. Explained that NFs are neural networks defining continuous functions over space, with examples: SDF/UDF for geometry and NeRF for appearance. Noted the advantage that memory footprint is decoupled from resolution, allowing arbitrarily fine detail from a fixed-size model. Conveyed the idea that the same MLP architecture can learn different field types, hinting at a unified representation for 3D objects.
Practical Applications of NFs: Outlined potential use cases and why NFs are exciting. For instance, generating and editing 3D models in gaming/VR via latent manipulation, using NFs in robotics for object manipulation and scene understanding, medical imaging uses (implicit representations of scans for analysis), and digital twins in industry stored as NFs. This slide set the stage that if we had a way to process NFs easily, it could impact many domains.
Key Terminology and Research Question: Defined implicit neural representation (INR) and clarified “object-centric” (focus on single objects rather than whole scenes). Then posed the central research question: Can we directly use neural field weights as input to deep learning pipelines? (E.g., classify an object from its NF without converting it to point cloud or image). This was identified as a non-trivial challenge since it means operating on the network parameters themselves rather than the explicit shape.
Prior Approaches (Brief): Briefly mentioned prior work for context. For example, Functa – a shared network approach that tried to embed multiple NFs into one model and the issues it faced fitting an entire dataset in one network. Mentioned that other recent methods (DWSNet, NFN, NFT) leverage weight space symmetries or specialized architectures for NFs, indicating an emerging research trend. This motivated the need for a new approach (nf2vec) focusing on individual NFs.
nf2vec Framework Overview: Introduced nf2vec as the proposed solution – “Neural Field to Vector” encoding. Visually, the slide (Fig. 1) showed how NFs (regardless of input modality) can be embedded into a latent code, which can then feed various tasks. Described the framework’s two parts: an encoder that processes NF weights and produces a compact embedding, and a decoder that can reconstruct the field’s outputs from the embedding. Emphasized that nf2vec looks only at the NF (its weights) to compute the embedding (no sampling of the function).
Encoder Architecture (Fig. 2): Detailed how the NF’s weights are organized for input. All weight matrices and biases are stacked into one matrix $P$ (with appropriate zero-padding for smaller layers). The encoder then applies a series of fully-connected layers with BatchNorm and ReLU to each row of $P$ independently. The slide likely included a schematic: showing $P$ going through layers (as individual row vectors) and then a final max-pooling across rows to produce the single embedding vector. Key point: this architecture scales to large NFs without a massive model, because of weight sharing across rows and pooling.
Weight Alignment & Initialization: Noted a critical implementation detail: to make row-wise processing meaningful, all NFs must share the same architecture and initial weight ordering. The slide mentioned that a shared random initialization across all neural fields was used to ensure weights are aligned. This helps the encoder learn consistent features from the same row index of $P$ across different objects. It was explained that without this, the encoder would struggle since one NF’s $i$-th row might not correspond to the same semantic part as another NF’s $i$-th row due to permutation symmetries. (The seminar likely highlighted this as a “critical consideration” for training nf2vec reliably.)
Decoder and Training Strategy: Described the decoder’s role: it takes the embedding and a coordinate to predict the field value. The training procedure mimics how the original NFs are trained. For shape fields (SDF/UDF/OF), they sample points in space and supervise the decoder with the known distance or occupancy values. For NeRF, they perform differentiable volume rendering and compare to ground-truth pixel colors of training images, using a weighted loss to focus on object pixels. The slide likely noted that the decoder predicts field values, not weights – emphasizing that we care about reconstructing the function’s output. Mentioned using techniques like positional encoding of coordinates (via a library like tiny-cuda-nn) to help the decoder represent high-frequency details, and skip connections to inject the embedding at multiple layers for better gradient flow (these implementation details were probably listed as bullet points).
Experimental Validation – Classification: Presented results showing nf2vec’s performance on shape classification benchmarks. Listed accuracy numbers: e.g., ModelNet40 ~87.0%, Manifold40 ~86.3%, ShapeNet10 ~93.0%. These were compared to baseline methods to illustrate nf2vec is on par with or close to specialized models. Highlighted that the same downstream classifier architecture was used for all representations, unlike baselines which differ for point clouds vs. meshes. Also showed a plot of inference time: nf2vec’s classification time stays roughly constant regardless of the NF’s internal resolution, whereas methods that require dense sampling slow down as resolution increases. Conclusion: nf2vec offers significant computational advantages by avoiding expensive pre-processing.
Qualitative Results – Latent Space Properties: Showed visualizations of what the embedding space has learned. For example, a slide with linear interpolation results: images of shapes morphing smoothly from one to another, and for NeRF, two objects’ appearances blending in a continuous transition. Another visualization was t-SNE clustering of embeddings: points in the plot colored by object class, demonstrating that categories form clusters naturally. Perhaps included were example reconstructions: decoding an embedding yields a shape nearly identical to the original NF’s shape, proving that the compression retained the necessary information. The takeaway was that the latent space is semantically structured and smooth, enabling meaningful interpolation and analogies.
Generality to All NF Types: This part underlined that nf2vec uses one encoder for all kinds of fields and one unified 1024-D embedding space. However, it requires different decoder heads depending on the NF type: a scalar output head for surface fields (SDF/UDF/OF) and a 4-channel (RGB+$\sigma$) output head for NeRF. The implementation details slide might have noted how the NeRF decoder includes special activation functions for RGB and density and normalizes coordinates within an axis-aligned bounding box. These details reassure that the framework was carefully engineered to handle both geometry-only and appearance+geometry cases.
Critical Implementation Considerations: A slide likely enumerated important training details:
Shared weight initialization (reiterated its necessity for convergence and meaningful embeddings).
Handling of SIREN networks: if sine activation NFs were used, the very first and last layers might not carry object-specific info, so those were ignored in encoding (as they “don’t contribute to downstream tasks”).
Data augmentation was done offline on the training set rather than online in each epoch, due to difficulty of augmenting NFs directly.
A concept of linear mode connectivity was mentioned: NFs trained with identical initialization exhibit a property that you can interpolate their weights without leaving the solution manifold. This was leveraged to argue that using the same init makes interpolation in weight space meaningful and helps nf2vec’s consistency. (This corresponds to an experiment where they showed interpolating between two NFs of the same shape yields low loss if they started from same init, indicating a flat minimum connecting them.)
Extension to Neural Radiance Fields: Discussed results specific to NeRF processing. nf2vec was applied to a dataset of object-centric NeRFs. It achieved around 87.3% accuracy in classifying NeRFs by object category, using the same simple MLP classifier as before. This is notable because the alternative (render multiple views and run image classifiers) is far slower. Indeed, nf2vec’s NeRF classification pipeline was cited to be two orders of magnitude faster than multi-view image-based methods. The presentation likely showed an example NeRF to NeRF generation: training a GAN to produce nf2vec codes for NeRFs resulted in plausible novel radiance fields (some images in Fig. 20 of the paper demonstrate generated views). Also shown was the cross-modal mapping from a shape embedding to a NeRF embedding: an input point cloud’s NF code mapped to an output NeRF code produced a colored 3D object (illustrated with a point cloud vs. rendered image side by side). These examples highlighted nf2vec’s capability to handle both geometry and appearance together.
Comparison with Recent Methods: A slide positioned nf2vec against other state-of-the-art NF processing approaches. It likely listed: DWSNet, NFN, NFT (all published in 2023-24) and Functa, DeepSDF (older approaches). Points made:
nf2vec outperforms DWSNet, NFN, NFT on classification accuracy of neural fields.
Methods that use a shared network for all data (like Functa, DeepSDF’s latent codes) cannot match the reconstruction quality of per-object networks – nf2vec and similar per-object methods maintain high fidelity.
nf2vec provides the best balance between accuracy and fidelity, whereas some others trade one for the other.
They introduced the first benchmark for NF classification, establishing a baseline for future research in this nascent area. This comparative analysis was to convince the audience that nf2vec is currently a leading technique in this niche.
Limitations and Challenges: The presentation transparently covered known issues:
Sometimes nf2vec’s accuracy is slightly below specialized networks designed explicitly for a certain 3D format (acknowledging the performance trade-off).
There is no straightforward way to augment data online when working with NF weights. This was highlighted as a limitation and an area for improvement.
In NeRF embeddings, geometry and appearance are entangled, limiting the ability to control them independently.
The requirement of shared initialization restricts flexibility – you can’t arbitrarily change the NF architecture or initialization without retraining the encoder. These points mirror the critical evaluation in this review, showing the authors were aware of and candid about these constraints.
Advanced Applications – Beyond Basics: A forward-looking or extended results section, showing additional cool things they did with nf2vec:
Shape Completion: Using a mapping in latent space to go from an embedding of an incomplete shape to a complete shape’s embedding, thereby performing completion without direct point cloud operations.
Cross-Representation Conversion: Demonstrating latent space mappings like UDF $\to$ SDF (point cloud to mesh conversion via code) or UDF $\to$ NeRF (adding texture to a shape). Indicated that simple MLP translators can be trained to convert between different NF modalities.
Latent Space Generative Modeling: Training a GAN (or diffusion model) on nf2vec codes for unconditional generation of new shapes. Showed that samples from a latent GAN decoded into realistic 3D geometries at arbitrary resolution. This hints at a data-driven way to generate 3D content using NFs.
These examples were likely presented to spark discussion on what can be done once we have a fixed-size vector representation for 3D objects – essentially opening the toolbox of deep learning for 3D (which traditionally lagged behind due to data format issues).
Performance & Scalability Analysis: Provided quantitative analysis reinforcing that nf2vec’s computational overhead is manageable:
Constant-Time Inference: The encoding time does not depend on how detailed the NF is (because even a high-resolution NF still has the same network architecture, just fitted better). So whether an NF represents a 16x16x16 volume or a 512x512x512 volume, the embedding time is the same — a big advantage over methods that scale with point or voxel count. A specific example: classification in 1.03 ms vs. 97.56 ms for a multi-view image approach.
Memory Usage: The memory required grows linearly with NF size, and the encoder’s row-wise processing means it doesn’t need to hold a massive flattened vector in memory all at once — it can stream through rows. Compared to an alternative of flattening and using a giant fully-connected layer (800M params in a naive encoder example), nf2vec’s encoder is far more compact.
Possibly a note that the method supports batching (so multiple NFs can be embedded concurrently on a GPU, further improving throughput).
This reassured that nf2vec can handle large datasets and high-complexity inputs at scale, which is important if we envision thousands of objects in a repository being processed.
Future Directions: Concluded the technical part with prospects for future research (these were explicitly enumerated in the talk):
Multi-Scale Processing: Incorporating hierarchical representations so that very large or complex objects can be represented at multiple levels of detail within the NF paradigm.
Disentangling Geometry and Appearance: Developing techniques to separate different factors in the embedding (e.g., have two codes or a structured code) to allow independent control, addressing the entanglement issue in NeRF.
Online Data Augmentation: Finding methods to directly perturb neural field representations (or their embeddings) in meaningful ways to generate new training examples without leaving the implicit domain.
Scene-level Understanding: Extending nf2vec from single objects to entire scenes, which would involve handling multiple objects’ NFs together or encoding spatial compositions. This might require new architectures or hierarchical embeddings.
Each of these points was likely briefly explained to inspire how the community could build on nf2vec to overcome its limitations or broaden its scope.
Conclusion and Impact: The final slide summarized the key takeaway: nf2vec demonstrates that direct processing of neural field representations is both feasible and effective. It provides a unified approach to 3D data, achieving competitive performance with much simpler pipelines. The method has clear practical advantages (speed, flexibility) that could make it useful in real-world settings. The presenters likely expressed that this work paves the way for NFs to become a standard in computer vision/graphics, bridging the gap such that one day 3D deep learning might not distinguish between different shape formats — all could be handled as neural fields. There may have even been a forward-looking remark (as per one slide’s rhetorical question about “Where could we be in 10 years?”) suggesting possibilities like large language models interacting with 3D neural fields, or neural fields becoming as common as images for AI tasks. The tone was optimistic, highlighting that while challenges remain, nf2vec is a step toward a more unified and powerful paradigm for 3D understanding.
(The outline above reflects the content and sequencing of the seminar slides “NeuralField.pptx”, capturing each major section and its focus.) \section{Attachments and Experiment Section Placeholder}
\[1ex] \textit{(No additional attachments are included in this review. A detailed experimental results section and any supplementary correspondence or appendices were omitted due to time constraints. In a full report, this is where extended experiment details and supporting materials would be provided.)} \bibliographystyle{acl_natbib}
\bibliography{custom}
\end{document}